\documentclass[11pt,a4paper]{article}

\usepackage[T1]{fontenc}
\usepackage[utf8]{inputenc}
\usepackage[spanish,es-noshorthands]{babel}
\usepackage{lmodern}
\usepackage{amsmath,amssymb,amsthm}
\usepackage[margin=2.4cm]{geometry}
\usepackage{booktabs}
\usepackage{microtype}
\usepackage{enumitem}
\usepackage{setspace}
\usepackage{titlesec}
\usepackage[colorlinks=true,linkcolor=black,citecolor=black,urlcolor=blue]{hyperref}

\theoremstyle{plain}
\newtheorem{teorema}{Teorema}
\newtheorem{corolario}[teorema]{Corolario}
\newtheorem{proposicion}[teorema]{Proposición}
\theoremstyle{definition}
\newtheorem{problema}{Problema abierto}
\theoremstyle{remark}
\newtheorem*{obs}{Observación}

\newcommand{\pc}{\ifmmode\text{\,\%}\else\,\%\fi}
\newcommand{\MER}{\mathrm{MER}}
\newcommand{\dd}{\,\mathrm{d}}
\newcommand{\p}{\partial}
\newcommand{\Rv}{R_v}
\newcommand{\F}{\mathcal{F}}
\newcommand{\R}{\mathbb{R}}

\setstretch{1.12}
\setlist[itemize]{leftmargin=1.3em,itemsep=0.25em,topsep=0.3em}
\setlist[enumerate]{leftmargin=1.4em,itemsep=0.25em,topsep=0.3em}
\titleformat{\section}{\large\bfseries}{\thesection.}{0.6em}{}
\titleformat{\subsection}{\normalsize\bfseries}{\thesubsection}{0.5em}{}

\title{\vspace{-1.3cm}
\textbf{¿Qué se puede recuperar de una erupción?}\\[0.35em]
\large Identificabilidad y no-unicidad en el ascenso bifásico\\
de magma por un conducto volcánico\\[0.7em]
\normalsize Propuesta de tesis de \textbf{doble titulación}\\
Geología \ + \ Ingeniería Matemática}
\author{Ignacio Gómez Gómez\\[0.3em]
{\small Universidad de Chile}\\
{\small Facultad de Ciencias Físicas y Matemáticas}\\
{\small Departamento de Geología \quad / \quad Departamento de Ingeniería Matemática}}
\date{\today}

\begin{document}
\maketitle
\vspace{-0.6cm}

\begin{abstract}\noindent
Esta propuesta formula una tesis que tiene que funcionar \emph{entera} para
las dos disciplinas: no es un trabajo de geología con un anexo numérico, ni
un trabajo de problemas inversos con un volcán de adorno. El objeto físico
es el ascenso de magma por un conducto (modelo bifásico 1D de Slezin y
Kozono--Koyaguchi, ya implementado en este repositorio). El objeto
matemático es el operador directo
$\F(\theta)=(q_{\mathrm{ex}}^\ast(\theta),\,q_{\mathrm{ef}}^\ast(\theta))$
y su inverso: dados observables de una erupción, qué se puede afirmar
sobre geometría, sobrepresión, volátiles y estilo, y con qué certeza.

La tesis \textbf{no} es ``hacer una inversión y reportar números''. Es
caracterizar \textbf{qué es recuperable y qué no}. Hay al menos dos
no-unicidades estructurales, independientes: (i)~geométrica, demostrada
en un recorte de lubricación (el $\MER$ ve $R(\cdot)$ sólo a través de
un funcional escalar $J_4$); (ii)~de estilo eruptivo, ya visible en el
operador bifásico (un mismo $\theta$ puede devolver solución explosiva,
efusiva, las dos, o ninguna). La pregunta central es si (i) sobrevive
al deslizamiento $u_g\neq u_m$, y qué datos restauran unicidad. El caso
principal es Calbuco 2015; el paper JGR 2025 aporta los demás casos
(Merapi, Caulle, Villarrica, St.\ Helens, \ldots) para el estilo.
\end{abstract}

\tableofcontents
\vspace{0.8em}

% ===========================================================================
\section{Por qué tiene que ser una sola tesis para las dos carreras}
\label{sec:doble}

Una tesis de doble titulación no se aprueba si cada profesor puede leer
sólo ``su'' mitad. Tiene que haber:

\begin{itemize}
\item una \textbf{pregunta geológica} que no se responde sin el aparato
  matemático (no basta simular Calbuco otra vez);
\item una \textbf{pregunta matemática} que no es un ejercicio abstracto
  (el operador, el dato y la no-unicidad vienen del volcán);
\item un \textbf{resultado} que ambas comunidades puedan citar: la
  geología, como advertencia sobre qué se puede inferir de una erupción;
  la matemática, como caracterización de un inverso mal puesto con
  dos fuentes de degeneración.
\end{itemize}

Esa pregunta, escrita para los dos, es:

\begin{quote}
\textit{Dados los observables de superficie de una erupción ---tasa de
descarga, estilo (explosivo / efusivo), masa, altura de columna, texturas---,
¿qué se puede recuperar del estado del magma en profundidad y de la
geometría del conducto, qué queda estructuralmente indeterminado, y qué
dato adicional reduce esa indeterminación?}
\end{quote}

\textbf{Lo que la geología necesita de la tesis.} No otro perfil
$(P,\phi,u)(z)$ hacia adelante. Una teoría de \emph{inferencia}:
cuándo el $\MER$ informa sobre el conducto profundo y cuándo no; cuándo
el estilo es un dato independiente y no un postprocesado; cómo
discriminar, con Calbuco, entre gatillo interno (segunda ebullición) y
gatillo por inyección. Eso cambia cómo se interpretan depósitos,
inclusiones y monitoreo.

\textbf{Lo que la ingeniería matemática necesita de la tesis.} Un
operador $\F$ bien definido (DAE de índice~1 + tiro), tres preguntas
clásicas de Hadamard (unicidad, estabilidad, reconstrucción), un
teorema de degeneración en un recorte, la pregunta abierta de si
sobrevive en el sistema completo, y un esquema de regularización /
inversión bayesiana con cuantificación de incertidumbre. Eso es un
problema inverso de verdad, no un ajuste de curvas.

Las secciones~\ref{sec:geo} y~\ref{sec:math} desarrollan cada lado.
La~\ref{sec:objetivos} separa objetivos por disciplina. La
estructura de capítulos (\S\ref{sec:capitulos}) está armada para que
ningún capítulo sea ``sólo de uno''.

% ===========================================================================
\section{Resumen ejecutivo}
\label{sec:resumen}

El modelo directo ya existe. Dado
$\theta=(R,\,\Delta P,\,c_0,\,T,\,\xi_0,\,\phi_{\mathrm{crit}},\ldots)$,
los códigos explosivo y efusivo
(\texttt{RIconduitex5\_5}, \texttt{RIconduitef5\_5})
resuelven \textbf{el mismo} DAE bifásico
$\mathbf{y}=(P,\phi,N,\xi,u_m,u_g)$
y buscan un caudal $q$ que cierre la condición de superficie. El
observable no es un escalar: es el par
\begin{equation}
\label{eq:Fpar}
\F(\theta)
=\bigl(q_{\mathrm{ex}}^\ast(\theta),\;q_{\mathrm{ef}}^\ast(\theta)\bigr)
\in
\bigl(\R_+\cup\{\varnothing\}\bigr)^2.
\end{equation}
Un mismo $\theta$ puede dar solo explosiva, solo efusiva, \textbf{las
dos}, o ninguna (Castruccio, Rebolledo \& Gómez, 2025). Invertir un
$\MER$ sin el estilo deja las dos ramas confundidas.

Tres resultados preliminares, obtenidos en un \emph{recorte} de
lubricación (\texttt{tesis/prototipo/}), definen el esqueleto
matemático. El recorte \textbf{no} reemplaza al bifásico: sirve para
demostrar.

\begin{enumerate}
\item \textbf{Degeneración geométrica $J_4$.}
  En el régimen viscoso dominado por fricción,
  $\MER\cdot J_4[R]=\mathrm{const}$, con
  $J_4[R]=\int_H^{z_f}R(z)^{-4}\,\dd z$.
  La preimagen de un dato es una hipersuperficie de codimensión~1 en
  un espacio de funciones. Verificado a precisión de máquina.
\item \textbf{La no-unicidad es enorme.}
  Seis geometrías con volúmenes que difieren en $+317\pc$ dan el mismo
  $\MER$ con error $10^{-7}$. El $90\pc$ de la resistencia está en los
  últimos metros: el conducto profundo es invisible para el caudal.
  \textbf{Esto está demostrado en el recorte.} La tesis pregunta si
  sobrevive en el bifásico.
\item \textbf{El tiempo separa escalares, no la función $R(\cdot)$.}
  Agregar $\MER(t)$ reduce el intervalo de $R$ constante de
  $[8,19]$~m a $[13,18]$~m, y deja $\Delta P$ con sólo una cota
  superior. Sobre $R(\cdot)$, cada punto de la curva
  $P_{\mathrm{cam}}\mapsto\MER_{ss}$ devuelve el mismo $J_4$.
\end{enumerate}

Hay una cuarta no-unicidad, ésta sí en el operador bifásico real:
\textbf{el estilo}. Merapi y Caulle son el ejemplo geológico: el
mismo sistema puede alimentar una columna y un domo.

La tesis recomendada, en una frase:

\begin{quote}
\textit{Identificabilidad y no-unicidad en modelos bifásicos de
ascenso de magma, con dos degeneraciones (geometría y estilo),
aplicada a Calbuco 2015 y a los casos del paper JGR 2025.}
\end{quote}

% ===========================================================================
\section{Planteamiento del problema}
\label{sec:planteamiento}

\subsection{El fenómeno}

Una erupción es el vaciado, total o parcial, de un reservorio de magma
a través de un conducto. En el conducto el magma descomprime: exsuelve
agua, nuclea y crece burbujas, puede cristalizar microlitos, cambia
órdenes de magnitud en viscosidad, y ---si la fracción de gas cruza un
umbral--- se fragmenta y pasa de flujo de líquido con burbujas a flujo
de gas con partículas. Eso decide si la erupción es \emph{explosiva}
(columna, tefra) o \emph{efusiva} (domo, lava). El estilo no es un
adjetivo: es una condición de borde distinta del mismo sistema.

Los observables de superficie ---tasa de descarga másica ($\MER$),
altura de columna, masa emitida, densidad y vesicularidad de
piroclastos, deformación, gas--- son funcionales de lo que ocurre
\emph{adentro}. La pregunta geológica de siempre (``¿cómo era el
conducto? ¿qué presión había en la cámara? ¿cuánta agua? ¿por qué
explotó?'') es, matemáticamente, un problema inverso.

\subsection{El operador}

Sea $\Theta$ el espacio de parámetros (puede contener funciones:
$R(\cdot)$, $\beta(\cdot)$, $P_{\mathrm{cam}}(\cdot)$) y $\mathcal{D}$
el de datos. El modelo directo define
\begin{equation}
\F:\Theta\longrightarrow\mathcal{D},\qquad
\theta\longmapsto d=\F(\theta).
\end{equation}
El dato observado es $d^{\mathrm{obs}}=\F(\theta^\dagger)+\varepsilon$.
Las tres preguntas de Hadamard:

\begin{enumerate}[label=(P\arabic*)]
\item \textbf{Unicidad / identificabilidad.}
  ¿Es $\F$ inyectivo? Si no, caracterizar $\F^{-1}(d)$.
\item \textbf{Estabilidad.}
  ¿Es $\F^{-1}$ continuo? ¿Con qué módulo?
\item \textbf{Reconstrucción.}
  Algoritmo, regularización, cuantificación de incertidumbre.
\end{enumerate}

En este problema la respuesta a (P1) es \textbf{no}, y se puede
demostrar en un recorte. Hay además no-unicidad de estilo en el
operador completo. Eso no es un defecto del código: es el contenido
de la tesis.

\subsection{Qué no es esta tesis}

\begin{itemize}
\item No es reescribir \texttt{RIconduitex5\_5.py} ni
  \texttt{RIconduitef5\_5.py}. Esos solvers \emph{son} el operador
  directo.
\item No es reemplazar el bifásico por un modelo de una sola
  velocidad. El recorte de \texttt{tesis/prototipo} es un
  laboratorio para teoremas.
\item No es ``correr MCMC hasta que salga un radio''. Un número sin
  barras de error, en un inverso degenerado, no significa nada.
\item No es un modelo 3D de conducto. $L/R\sim 400$ en Calbuco; la
  axisimetría captura la física. El 3D no paga frente al costo, y
  el inverso se vuelve intratable.
\end{itemize}

% ===========================================================================
\section{El lado geológico}
\label{sec:geo}

\subsection{Por qué el conducto importa}

El conducto es el filtro entre la cámara y la superficie. Su radio,
su rugosidad, su posible ensanchamiento sineruptivo y el estado del
magma que lo recorre ($c_0$, $T$, $\xi$, $\phi$) controlan el
$\MER$ y el estilo. Dos erupciones con la misma cámara pueden verse
muy distintas si cambia el conducto; dos conductos distintos pueden
verse iguales en superficie (ésa es la degeneración $J_4$).

La geología de terreno lee el conducto \emph{después}: depósitos,
clastos líticos, texturas de borde frente a centro, geometría del
cráter, diques aflorantes. El modelo de conducto lee el conducto
\emph{durante}. El problema inverso es el puente entre las dos
lecturas.

\subsection{Estilo eruptivo: un dato, no un postprocesado}

La distinción explosivo / efusivo se ha tratado a menudo como un
umbral empírico (contenido de agua, viscosidad, $\MER$). El modelo
bifásico la convierte en un hecho matemático: el mismo DAE admite
dos cierres de superficie.

\begin{center}
\small
\begin{tabular}{@{}lll@{}}
\toprule
 & explosivo (\texttt{ex}) & efusivo (\texttt{ef}) \\
\midrule
$P(0)$ & $P_{\mathrm{atm}}$ \textbf{o} choque $u_g\approx\sqrt{\Rv T}$
  & $P_{\mathrm{atm}}$ \\
$\phi(0)$ & $\phi\ge\phi_{\mathrm{crit}}$ (fragmentó)
  & $\phi\le\phi_{\mathrm{crit}}$ (no fragmentó) \\
$z_{\mathrm{exit}}$ & $\ge -5\,\mathrm{m}$ & $\in(-2,0]$ \\
búsqueda de $v_{\mathrm{in}}$ & $[0.1,50]$, guess $10\,\mathrm{m\,s^{-1}}$
  & $[10^{-5},5]$, guess $0.1\,\mathrm{m\,s^{-1}}$ \\
\bottomrule
\end{tabular}
\end{center}

\vspace{0.3em}
Para un mismo $\theta$ el par~\eqref{eq:Fpar} puede ser
$(q,0)$, $(0,q)$, $(q_{\mathrm{ex}},q_{\mathrm{ef}})$ o
$(\varnothing,\varnothing)$. Eso ya está en el paper
(Castruccio \emph{et al.}, 2025) y en el notebook
\texttt{mainconduit\_ambos.ipynb}: Merapi y Caulle son los casos
donde las dos ramas conviven. \textbf{El dato observado no es
sólo un $\MER$. Es el par $(\MER,\,\mathrm{estilo})$.}

Consecuencia para el inverso: si se invierte sólo el caudal, las
dos ramas se confunden. Si se invierte el par, se identifica la
rama, pero $R$ y $\Delta P$ siguen sin ser únicos ---en el paper
ni siquiera son datos: los calcula un cierre de estabilidad del
conducto.

\subsection{Calbuco, abril de 2015}

Caso principal. Erupción sub-Pliniana en dos fases, andesita
basáltica, muy bien documentada, chilena, y ya parametrizada en
este repositorio (\texttt{calbuco2015d.py}).

\begin{center}
\footnotesize
\begin{tabular}{@{}p{0.28\textwidth}p{0.38\textwidth}p{0.26\textwidth}@{}}
\toprule
observable & valor & fuente \\
\midrule
volumen total (no DRE) & $0.27$--$0.56\,\mathrm{km}^3$
  & Romero; Castruccio; Van Eaton; Pardini \\
reparto de fases & $\sim 38\pc$ / $\sim 62\pc$
  & Castruccio et al.\ (2016) \\
masa fase 1 / 2 & $8\times10^{10}\,\mathrm{kg}$ en $1.5\,\mathrm{h}$ /
  $3.2\times10^{11}\,\mathrm{kg}$ en $6\,\mathrm{h}$
  & ídem \\
$\MER$ media & $1.4\times10^7$ y $1.5\times10^7\,\mathrm{kg\,s^{-1}}$
  & ídem \\
$\MER$ máx.\ (isopletas) & $2.4$ y $2.7\times10^7\,\mathrm{kg\,s^{-1}}$
  & ídem \\
altura de columna & $15$--$17\,\mathrm{km}$ sobre el cráter
  & SERNAGEOMIN; Van Eaton; Romero \\
composición & $54$--$55\pc$ SiO$_2$ (depósitos);
  $62.72\pc$ en el input del modelo
  & Romero; \texttt{calbuco2015d} \\
almacenamiento & $8$--$12\,\mathrm{km}$, $900$--$950\,^{\circ}\mathrm{C}$,
  $5.5$--$6.5\,\mathrm{wt\pc}$ H$_2$O
  & Arzilli (2019); Morgado (2019) \\
deformación & co-eruptiva detectada; pre-eruptiva bajo ruido
  & Delgado et al.\ (2017) \\
pluma y rayos & serie temporal
  & Van Eaton et al.\ (2016) \\
\bottomrule
\end{tabular}
\end{center}

Dos hipótesis \emph{publicadas y contrapuestas} sobre el gatillo:

\begin{itemize}
\item \textbf{Arzilli et al.\ (2019):} gatillo interno. Cristalización
  prolongada $\to$ segunda ebullición $\to$ sobrepresurización, sin
  intrusión de magma fresco.
\item \textbf{Morgado et al.\ (2019):} calentamiento localizado por
  inyección de magma caliente.
\end{itemize}

Ambas predicen historias $P_{\mathrm{cam}}(t)$ distintas y por tanto
curvas $\MER(t)$ distintas. Discriminarlas con factores de Bayes es
un resultado geológico concreto, del tipo que un inverso bien
planteado puede responder. Van Eaton et al.\ (2016) aportan la serie
temporal más usable (altura de pluma y actividad eléctrica).

\textbf{Inconsistencia que hay que cerrar con el profesor de
geología.} El input del código usa $T=970\,^{\circ}\mathrm{C}$,
$c_0=4\,\mathrm{wt\pc}$, $H=-7\,\mathrm{km}$. La petrología reciente
sugiere $900$--$950\,^{\circ}\mathrm{C}$, $5.5$--$6.5\,\mathrm{wt\pc}$
y $8$--$12\,\mathrm{km}$. La viscosidad de Giordano es
exponencialmente sensible a $T$ y $c_d$. Las previas de la inversión
no se pueden fijar sin resolver esto.

\subsection{Los otros casos (paper JGR 2025)}

El paper (Castruccio, Rebolledo \& Gómez, 2025) barre composición,
$T$, cristales y agua, y muestra que el estilo y el $\MER$ salen del
mismo modelo. Los \texttt{.m} están en \texttt{casos/}, sin editar.
Sirven para la no-unicidad de estilo y como banco de prueba del
operador~\eqref{eq:Fpar}.

\begin{center}
\footnotesize
\begin{tabular}{@{}p{0.22\textwidth}p{0.22\textwidth}p{0.32\textwidth}c@{}}
\toprule
erupción & estilo observado & por qué entra & geom. \\
\midrule
Calbuco 2015 & sub-Pliniana & caso inverso principal & cil. \\
Vesuvius 79 & Pliniana & extremo explosivo, $c_0$ alto & cil. \\
Huaynaputina 1600 & Pliniana & ídem, riolita & cil. \\
Caulle 2011 & sub-Pliniana $\to$ lava & \textbf{las dos ramas} & cil. \\
Merapi 2010 & explosiva + domo & \textbf{las dos ramas} & cil. \\
St.\ Helens 2004 & domo & extremo efusivo (Anderson--Segall) & cil. \\
Villarrica 2015 & hawaiana / estromboliana & Henry basáltico, dique & dique \\
Pinatubo, Quizapu & Plinianas & extras, $R$ grande & cil. \\
\bottomrule
\end{tabular}
\end{center}

Villarrica usa $C_1=6.80\times10^{-8}$, $\beta=0.7$ y geometría de
dique. El resto, Henry silíceo y cilindro. En Vesuvius y
Huaynaputina el \texttt{.m} deja $R$ y $\Delta P$ comentados: el
catálogo usa $16\,\mathrm{m}$ / $5\,\mathrm{MPa}$ como partida, no
como dato.

\subsection{Qué aportaría esto a la volcanología}

\begin{enumerate}
\item Una \textbf{advertencia metodológica demostrable}: inferir el
  conducto profundo desde el $\MER$ es, en el recorte, inferir la
  previa. Si eso sobrevive en el bifásico, cambia cómo se leen
  radios ``invertidos'' en la literatura.
\item Tratar el \textbf{estilo como dato} del inverso, no como
  etiqueta a posteriori. Eso conecta el paper 2025 con la tesis.
\item El primer intento de inversión bayesiana con modelo físico de
  una erupción \textbf{explosiva} (la línea Anderson--Segall / Wong
  es toda efusiva: St.\ Helens 2004--2008).
\item Un caso chileno, con implicancias de amenaza (columnas de
  $15$--$17\,\mathrm{km}$, aeropuertos, lahares). El diseño óptimo
  de observables (qué medir para reducir más la posterior) es
  directamente útil para monitoreo.
\item Discriminar gatillo interno vs.\ inyección en Calbuco, si los
  datos temporales alcanzan.
\end{enumerate}

\subsection{Lo que la geología \emph{no} le pide a esta tesis}

No le pide un mapa 3D del sistema de alimentación, ni un modelo de
columna eruptiva, ni rehacer la petrología de Calbuco. Le pide
decir, con el modelo de conducto que ya usamos, \emph{qué parte de
la historia profunda queda determinada por los depósitos y el
monitoreo, y qué parte no}.

% ===========================================================================
\section{El lado de ingeniería matemática}
\label{sec:math}

\subsection{Tipo de problema}

Es un problema inverso para un DAE de índice~1 con un parámetro de
tiro, y ---si se deja $R(\cdot)$ libre--- un problema inverso en
dimensión infinita. El dato es de frontera (superficie). La
incógnita puede ser un escalar ($\Delta P$, $c_0$), un par
$(R,\Delta P)$, o una función $R(\cdot)$.

Formulación determinista:
\begin{equation}
\label{eq:tikhonov}
\theta_\alpha\in\arg\min_{\theta\in\Theta}
\tfrac12\bigl\|\F(\theta)-d^{\mathrm{obs}}\bigr\|_{\Sigma^{-1}}^2
+\alpha\,\mathcal{R}(\theta),
\end{equation}
con $\mathcal{R}$ de Tikhonov, de variación total (geometrías tipo
NC2 de Aravena) o de Sobolev. Elección de $\alpha$ por Morozov o
GCV.

Formulación bayesiana (la adecuada aquí):
\begin{equation}
\label{eq:bayes}
\pi(\theta\mid d^{\mathrm{obs}})
\propto
\exp\Bigl(-\tfrac12\bigl\|\F(\theta)-d^{\mathrm{obs}}\bigr\|_{\Sigma^{-1}}^2\Bigr)
\cdot\pi_0(\theta).
\end{equation}
Razones geológicas concretas: la petrología entrega previas reales
(inclusiones $\to c_0$; fases $\to T$ y profundidad; CSD $\to\xi$);
el $\MER$ se conoce típicamente salvo un factor~$2$, asimétrico.
Una estimación puntual sin posterior no sirve. Es el camino de
Anderson \& Segall (2013) y Wong et al.\ (2017, 2020), todavía no
aplicado a un caso explosivo ni a $R(\cdot)$ como función.

\subsection{El operador directo, escrito}

El estado en cada $z\in[H,0]$ es
\begin{equation}
\mathbf{y}(z)=\bigl(P,\;\phi,\;N,\;\xi,\;u_m,\;u_g\bigr).
\end{equation}
Las velocidades son algebraicas (masa + $q$ constante). $P$,
$\phi$, $N$ y $\xi$ son diferenciales. El sistema es
\begin{equation}
M\mathbf{y}'=F(\mathbf{y};q),\qquad
M=\mathrm{diag}(1,1,1,1,0,0).
\end{equation}
Masa:
\begin{equation}
q=\rho_m(1-\phi)u_m+\rho_g\phi u_g,\qquad
\rho_g=P/(\Rv T),\qquad
n=\frac{\rho_g\phi u_g}{q}.
\end{equation}
Momentum de cada fase (Kozono \& Koyaguchi, 2009):
\begin{align}
\rho_m(1-\phi)u_m u_m'
&=-(1-\phi)P'-\rho_m(1-\phi)g+F_{mg}-F_{mw},\\
\rho_g\phi u_g u_g'
&=-\phi P'-\rho_g\phi g-F_{mg}-F_{gw}.
\end{align}
Sustituir $u_m(P,\phi,\xi;q)$ y $u_g$ da un $2\times 2$ para
$(P',\phi')$. Solubilidad de Henry $c_s(P)=C_1 P^\beta$; cinética
de cristales de primer orden; coalescencia tipo Slezin; viscosidad
Giordano $\times$ \texttt{fvrel} $\times$ Llewellin; cuatro
regímenes de arrastre. La formulación completa está en
la nota \texttt{matematica\_modelo\_bifasico}. Eso \textbf{es}
el operador. No se recorta para invertir: se recorta para
demostrar.

El caudal $q$ no es dato: es el parámetro de un problema de tiro.
Hay dos cierres, \S\ref{sec:geo}. Éxito numérico: el contador de
iteraciones $<49$.

\subsection{Tres fuentes de no-unicidad (hay que no mezclarlas)}

\begin{enumerate}
\item \textbf{Geométrica ($J_4$).}
  Demostrada en el recorte $u_m=u_g$. Pregunta abierta en el
  bifásico. Es no-inyectividad de $\F$ respecto de $R(\cdot)$ a
  dato fijo.
\item \textbf{De estilo.}
  Ya en el bifásico: $\F(\theta)$ es un par, no un escalar. Misma
  $\theta$, dos $q$ posibles. Distinta de (1).
\item \textbf{Multiplicidad de estados estacionarios.}
  Curva sigmoidal $p_{\mathrm{ch}}$--$q$ (Melnik \& Sparks 1999;
  Kozono \& Koyaguchi 2012): a \emph{parámetros fijos}, varios $q$.
  Es no-unicidad del \emph{directo}, no del inverso. Complementaria.
\end{enumerate}

Anderson \& Segall (2013) ya notan que sus datos determinan
$R^4/8\eta$, no $R$ y $\eta$ por separado. Lo que no está hecho:
tratar $R(\cdot)$ como función, caracterizar el espacio nulo,
preguntar si sobrevive al slip, y tratar el estilo como segunda
degeneración.

\subsection{El recorte, y el teorema}

Para demostrar (P1) hace falta separar variables. Se recorta el
bifásico ---no se cambia el volcán---: $u_m=u_g$, $\xi$ fijo,
exsolución en equilibrio, dominio hasta $z_f$. Queda
\begin{equation}
\label{eq:lub}
\frac{\dd P}{\dd z}
=-\rho(P)\,g-\frac{8\mu(P)\,\MER}{\pi\rho(P)\,R(z)^4},
\end{equation}
con $P(H)=P_{\mathrm{cam}}$ y $P(z_f)=P_{\mathrm{frag}}$.

\begin{teorema}[Degeneración geométrica]
En el régimen dominado por fricción ($\rho g$ despreciable),
\begin{equation}
\label{eq:J4}
\MER\cdot J_4[R]
=\frac{\pi}{8}\Phi(P_{\mathrm{cam}},P_{\mathrm{frag}}),\qquad
J_4[R]=\int_H^{z_f}R(z)^{-4}\,\dd z,\qquad
\Phi=\int_{P_{\mathrm{frag}}}^{P_{\mathrm{cam}}}\frac{\rho}{\mu}\,\dd P.
\end{equation}
El $\MER$ depende de $R(\cdot)$ \textbf{únicamente} a través de
$J_4[R]$.
\end{teorema}

\begin{proof}
Despreciar $\rho g$ en~\eqref{eq:lub} y separar variables. El lado
izquierdo no depende de $R(\cdot)$.
\end{proof}

\begin{corolario}
$\F_{\mathrm{recorte}}$ no es inyectivo en $R(\cdot)$. El conjunto
de nivel $\{R:J_4[R]=c\}$ es de dimensión infinita.
\end{corolario}

\begin{corolario}[Espacio nulo]
$DJ_4[R]\,\delta R=-4\int R^{-5}\delta R\,\dd z$, luego
\[
\mathcal{N}_R
=\Bigl\{\delta R\in C([H,z_f]):
\int_H^{z_f}R^{-5}\delta R\,\dd z=0\Bigr\}
\]
es un subespacio cerrado de \textbf{codimensión~1}.
\end{corolario}

El resultado no es un artefacto de Poiseuille. Si la fricción
factoriza $P'=-F(P)\,f(\MER)\,R^{-n}$, el dato ve sólo
$J_n[R]=\int R^{-n}\dd z$ ($n=4$ laminar, $19/4$ Blasius, $5$
Darcy--Weisbach).

\begin{proposicion}
Bajo las hipótesis del teorema, la curva
$P_{\mathrm{cam}}\mapsto\MER_{ss}$ determina $J_4[R]$ y nada más
de $R(\cdot)$.
\end{proposicion}

\begin{obs}
El arrastre $F_{mg}$ y el slip $u_g-u_m$ \emph{podrían} romper la
degeneración: introducen el radio de burbuja $r_b(\phi,N)$, una
escala que el recorte no tiene. Ésa es la pregunta matemática
central, no un resultado cerrado.
\end{obs}

\subsection{Identificabilidad práctica}

Aunque $\F$ no sea inyectivo, algunos parámetros pueden estar
prácticamente identificados. Herramientas:

\begin{itemize}
\item matriz de sensibilidad / Jacobiano
  $D\F[\theta^\dagger]$, SVD, valores singulares
  (grado de mal planteamiento);
\item matriz de información de Fisher
  $I(\theta)=D\F^{\mathsf T}\Sigma^{-1}D\F$;
\item perfiles de verosimilitud (Raue et al., 2009): para cada
  coordenada $\theta_i$,
  $\min\{\,{-}2\log L:\theta_i=\mathrm{fijo}\,\}$
  y umbral $\chi^2_{1,0.95}$.
\end{itemize}

En el recorte, en el punto $(R,\Delta P)=(16\,\mathrm{m},\,5\,\mathrm{MPa})$,
\[
J
=\frac{\p(\log\MER,\,\log G)}{\p(\log R,\,\log\Delta P)}
=\begin{pmatrix}4.000&0.195\\4.000&0.115\end{pmatrix},
\quad
G=\frac{\dd\MER}{\dd P_{\mathrm{cam}}}.
\]
El $4.000$ es Poiseuille. El $0.195$ es chico porque $\Delta P$ es
una fracción del empuje (el resto es flotabilidad): la compensación
local es $\Delta P\sim R^{-20.5}$. Con ruido realista
$\sigma_{\log\MER}=\ln 2/2$,

\begin{center}
\small
\begin{tabular}{@{}lcc@{}}
\toprule
datos & IC $95\pc$ de $R$ & IC $95\pc$ de $\Delta P$ \\
\midrule
sólo estacionario & $[8,19]\,\mathrm{m}$ (abierto)
  & $[0.2,250]\,\mathrm{MPa}$ (todo el rango) \\
estacionario + $\MER(t)$ & $[13,18]\,\mathrm{m}$
  & $[0.25,44]\,\mathrm{MPa}$ (sólo cota superior) \\
\bottomrule
\end{tabular}
\end{center}

Éste es un resultado, no un fracaso: el tiempo identifica el radio
escalar y acota $\Delta P$ por arriba. Para acotarla por abajo hacen
falta deformación o gas.

\subsection{Cuándo el término transiente importa}

Hay que no confundir ``el dato depende del tiempo'' con ``el
conducto está fuera de equilibrio''. El criterio es un Deborah
$\mathrm{De}=\tau_{\mathrm{conducto}}/\tau_{\mathrm{forzamiento}}$.
En el recorte, medido: $\tau_{\mathrm{hidr}}=8\,\mathrm{s}$,
$\tau_{\mathrm{resid}}=6\,\mathrm{min}$,
$\tau_{\mathrm{cam}}=5.83\,\mathrm{h}$ ($V=50\,\mathrm{km}^3$),
luego $\mathrm{De}=5.7\times10^{-4}$. El error del
cuasi-estacionario es $0.08\pc$ en $\MER(t)$ para Calbuco. Cruza
$1\pc$ a $\mathrm{De}\approx 0.09$ (forzamientos $\lesssim 90\,\mathrm{s}$).

Pero el $\tau$ que de verdad puede volver transiente una erupción
no es el hidráulico: es exsolución en desequilibrio
($10^1$--$10^3\,\mathrm{s}$), cristalización de microlitos
($10^4$--$10^6\,\mathrm{s}$), escape de gas
($10^3$--$10^5\,\mathrm{s}$). Con $\tau\sim10^4\,\mathrm{s}$, una
fase de $6\,\mathrm{h}$ pasa a $\mathrm{De}\approx 0.5$.

\textbf{Consecuencia:} lo invertible depende del régimen. En
cuasi-estacionario los datos informan geometría y cámara, y
\textbf{nada} sobre cinética. Sólo con $\mathrm{De}\gtrsim 0.1$ las
constantes de tasa son observables. Calbuco está cómodo en el
primero si se mira la hidráulica, y en la frontera si se mira la
cristalización. Un domo cíclico está en el segundo \emph{por
construcción} (el período lo fija la cinética).

\subsection{Una joya opcional: Webster y el espectro}

Linealizando flujo compresible en un tubo de sección $A(z)$ se
obtiene la ecuación de Webster, que con $p=A^{-1/2}\psi$ es
Sturm--Liouville con potencial
$V=(A^{1/2})''/A^{1/2}$. El teorema de dos espectros de Borg
(1946): \textbf{una sola familia espectral no determina $V$}. El
espectro de resonancia (LP, tremor) no determina la geometría.
Aun conociendo $V$, $(A^{1/2})''=VA^{1/2}$ tiene un espacio de
dimensión~2: hay una familia biparamétrica de secciones
acústicamente indistinguibles. Como en volcanología se infieren
radios desde LP, exhibir esa familia con parámetros realistas es
una advertencia fuerte y matemáticamente limpia. Puede ser
capítulo o sección, según el profesor de matemáticas.

\subsection{Lo que la matemática \emph{no} le pide a esta tesis}

No le pide existencia y unicidad del DAE directo en un espacio de
Sobolev abstracto (el código ya integra con IDA). Le pide
\emph{el inverso}: inyectividad, módulo de estabilidad,
regularización, y un algoritmo que respete la degeneración en vez
de esconderla.

% ===========================================================================
\section{Qué hay construido}
\label{sec:construido}

\begin{center}
\small
\begin{tabular}{@{}p{0.28\textwidth}p{0.38\textwidth}p{0.26\textwidth}@{}}
\toprule
módulo & formulación & estado \\
\midrule
\texttt{RIconduitex5\_5.py} &
  DAE 6 variables, tiro explosivo &
  funciona; IDA (\texttt{scikits.odes}) \\
\texttt{RIconduitef5\_5.py} &
  \textbf{el mismo DAE}, tiro efusivo &
  funciona; el par $(q_{\mathrm{ex}},q_{\mathrm{ef}})$ \\
\texttt{mainconduit\_ambos.ipynb} &
  tira las dos ramas, 6 figuras encima &
  listo; corre donde ya corre \texttt{mainconduit5\_5} \\
\texttt{RIconduit1D\_transient.py} &
  transporte $(\phi,N,\xi)$ + $P$ cuasi-estática; IMEX &
  funciona; sólo numpy \\
\texttt{RIconduit2D\_FD.py} &
  $u_m(r)$ por FD, Thomas + Picard &
  funciona; sólo numpy \\
\texttt{casos/} + \texttt{catalogo.py} &
  9 erupciones del paper 2025 &
  inputs sin editar \\
\texttt{tesis/prototipo/} &
  recorte de lubricación, exp.\ 1--4 &
  teorema $J_4$ + identificabilidad \\
\bottomrule
\end{tabular}
\end{center}

Tres observaciones que hay que poner sobre la mesa:

\begin{itemize}
\item \textbf{Los tres solvers grandes están desacoplados.}
  El 2D no usa el transiente. Un inverso serio necesita
  \textbf{un} operador. El de la tesis es el DAE bifásico; el
  transiente y el 2D son variantes, no otro proyecto.
\item \textbf{$\F$ no es diferenciable.}
  Hay $\max(0,\cdot)$, cambios de régimen, criterios de choque,
  bisecciones. Eso descarta gradiente hasta regularizar la
  fragmentación.
\item \textbf{Inconsistencias entre versiones}
  ($\phi_{\mathrm{crit}}$ por capilaridad en el estacionario,
  $0.75$ fijo en el transiente; permeabilidad definida y no
  acoplada). Antes de invertir hay que fijar \textbf{una} física.
\end{itemize}

Ninguna es un defecto del trabajo previo: son el primer capítulo
natural (convertir un código de simulación en un código de
inversión).

% ===========================================================================
\section{Estado del arte y nicho}
\label{sec:arte}

Revisión anotada en \texttt{tesis/ESTADO\_DEL\_ARTE.md}. Aquí el
nicho, en una tabla.

\begin{center}
\small
\begin{tabular}{@{}p{0.22\textwidth}p{0.34\textwidth}p{0.36\textwidth}@{}}
\toprule
dimensión & qué hay & qué falta \\
\midrule
tiempo &
  Melnik \& Sparks; La Spina et al.; Wong \& Segall 2019 &
  usarlo \textbf{para invertir}, no sólo simular; cuantificar
  cuánta información aporta \\
radial (2D) &
  Massol \& Jaupart; Collier \& Neuberg &
  nadie lo usó en un inverso; la pared como incógnita (Robin)
  está virgen \\
3D &
  escaso; $L/R$ enorme &
  no recomendado \\
geometría $R(z)$ &
  DOMEFLOW; Aravena; Costa (estudio \emph{directo}) &
  nadie la trató como \textbf{incógnita}; nadie caracterizó el
  nulo \\
inverso bayesiano &
  Anderson \& Segall 2013; Wong et al.\ 2017, 2020
  (todos \textbf{efusivos}) &
  ningún caso explosivo sub-Pliniano; ningún análisis
  estructural de no-unicidad; ningún caso chileno;
  el estilo no entra como dato \\
\bottomrule
\end{tabular}
\end{center}

\vspace{0.4em}
\textbf{El nicho es la intersección:} geometría (y estilo) como
incógnitas de un inverso, degeneración demostrada, datos que la
reducen, erupción explosiva chilena. Álvaro Aravena trabaja la
geometría no cilíndrica en el directo: contacto natural y posible
evaluador.

La no-unicidad que documenta la literatura clásica (Slezin 2003;
Melnik 2000; Kozono \& Koyaguchi 2012) es respecto de la tasa a
parámetros fijos. La de esta tesis es la otra flecha: respecto de
los parámetros y de $R(\cdot)$ a dato fijo.

% ===========================================================================
\section{Pregunta de investigación e hipótesis}
\label{sec:pregunta}

\paragraph{Pregunta principal.}
¿Es inyectivo el operador bifásico $\F$ respecto de la geometría
$R(\cdot)$ y respecto del estilo, y qué observables restauran
unicidad o, al menos, identificabilidad práctica?

\paragraph{Preguntas derivadas.}
\begin{enumerate}
\item ¿Sobrevive $J_4$ (o un $J_n$) cuando $u_g\neq u_m$?
\item El estilo, ¿es un funcional independiente de $R(\cdot)$, o
  está acoplado a la misma degeneración?
\item ¿Cuánto reduce $\MER(t)$ el nulo, comparado con un $\MER$
  escalar? ¿Y un perfil de porosidad de piroclastos? ¿Y
  deformación?
\item En Calbuco 2015, ¿qué cota sobre $\Delta P$ imponen los
  datos, y qué hipótesis de gatillo es más verosímil?
\item ¿En qué régimen (Deborah) las constantes de tasa son
  observables?
\end{enumerate}

\paragraph{Hipótesis de trabajo.}
\begin{enumerate}
\item En el bifásico la degeneración geométrica \emph{sobrevive},
  quizá con un funcional distinto de $J_4$ (el slip introduce
  $r_b$, pero la fricción de pared sigue factorizando en
  $R^{-2}$). El nulo puede dejar de ser de codimensión~1, pero no
  se espera unicidad.
\item El tiempo no recupera $R(\cdot)$; sí separa $(R,\Delta P)$
  cuando $R$ es constante.
\item El estilo es una no-unicidad \emph{ortogonal} a $J_4$: se
  levanta observando si fragmentó, no midiendo mejor el $\MER$.
\item Calbuco, hidráulicamente, es cuasi-estacionario: los datos
  no identifican $\tau_c$. Si se quiere invertir cinética, el
  caso tiene que ser un domo cíclico (St.\ Helens 2004, Merapi).
\end{enumerate}

% ===========================================================================
\section{Objetivos}
\label{sec:objetivos}

\subsection{Objetivo general}

Caracterizar la identificabilidad del modelo bifásico de conducto
---qué parámetros y qué funciones son recuperables desde
observables de erupción, y cuáles no--- y aplicarlo a Calbuco 2015
y a los casos del paper 2025, de modo que el resultado sea a la
vez un teorema (o un contraejemplo numérico preciso) y una
interpretación volcanológica.

\subsection{Objetivos específicos de Geología}

\begin{enumerate}[label=G\arabic*.]
\item Fijar, con evidencia petrológica, las previas de
  $\theta$ para Calbuco (reconciliar $T$, $c_0$, $H$ del código
  con Arzilli / Morgado) y el vector de datos
  $d^{\mathrm{obs}}$ (MER por fases, masa, columna, texturas,
  deformación, pluma).
\item Interpretar el par $(q_{\mathrm{ex}},q_{\mathrm{ef}})$ como
  dato de estilo, y mapear en los casos del paper 2025 cuándo el
  modelo predice las dos ramas (Caulle, Merapi) de acuerdo con lo
  observado.
\item Decir qué se puede afirmar, y qué no, sobre el conducto
  profundo de Calbuco a partir del $\MER$; en particular, una
  cota sobre $\Delta P$ y una discusión honesta de $R(\cdot)$.
\item Si los datos temporales alcanzan, comparar las dos
  hipótesis de gatillo por factores de Bayes.
\item Traducir el resultado a una recomendación de monitoreo:
  qué observable adicional reduce más la posterior (diseño
  óptimo, en versión geológica).
\end{enumerate}

\subsection{Objetivos específicos de Ingeniería Matemática}

\begin{enumerate}[label=M\arabic*.]
\item Escribir $\F$ como operador entre espacios (escalares y,
  para $R(\cdot)$, un Sobolev o $BV$): dominio, dato, regularidad
  a trozos por los cambios de régimen.
\item Demostrar el teorema $J_4$ y la caracterización del nulo en
  el recorte (hecho); estudiar el decaimiento de valores
  singulares de $D\F$ (grado de mal planteamiento).
\item Responder, numérica o analíticamente, si la degeneración
  sobrevive en el DAE bifásico. Entregable: contraejemplo
  (dos $R(\cdot)$ con el mismo par $(q_{\mathrm{ex}},q_{\mathrm{ef}})$)
  o enunciado de unicidad condicional.
\item Construir la tabla de identificabilidad (Fisher + perfiles
  de verosimilitud) de los parámetros finitos frente a cada
  conjunto de datos: estacionario, $\MER(t)$, estilo, texturas.
\item Plantear y resolver~\eqref{eq:tikhonov} o~\eqref{eq:bayes}
  para Calbuco, con regularización que respete el nulo (no
  pretender recuperar lo que $J_4$ esconde).
\item (Si el profesor de matemáticas lo pide.) Familia de Webster
  / Borg, o el adjunto del transiente, o el coeficiente de Robin
  en la pared.
\end{enumerate}

\subsection{Objetivos compartidos (los que amarran la doble titulación)}

\begin{enumerate}[label=C\arabic*.]
\item Un solo operador directo, el bifásico, con las dos ramas de
  tiro, y una decisión explícita sobre transiente y 2D.
\item Un capítulo de no-unicidad que un geólogo pueda citar
  (``el conducto profundo no se ve'') y un matemático pueda
  verificar (el funcional, el nulo, el experimento).
\item Un capítulo de Calbuco que un matemático pueda leer como
  inversión con incertidumbre, y un geólogo como interpretación
  del sistema de 2015.
\end{enumerate}

% ===========================================================================
\section{Metodología}
\label{sec:metodo}

\subsection{El operador no se reescribe}

Se llama a \texttt{RIconduitex5\_5\_f} y \texttt{RIconduitef5\_5\_f}.
El notebook \texttt{mainconduit\_ambos.ipynb} es el laboratorio de
las dos ramas. El recorte queda para teoremas y para explorar el
inverso a bajo costo (miles de evaluaciones, sin IDA).

\subsection{Etapa 0 --- Un operador invertible}

\begin{itemize}
\item Fijar $\phi_{\mathrm{crit}}$, permeabilidad, cinética: una
  física.
\item Regularizar la fragmentación (transición suave en $\phi$)
  para que $\F$ sea al menos Lipschitz, idealmente
  diferenciable a trozos.
\item Tests de regresión contra los solvers actuales.
\item Decidir si el transiente se usa como operador o sólo como
  generador de $\MER(t)$ cuasi-estacionario (el Deborah de
  Calbuco sugiere lo segundo para hidráulica).
\end{itemize}

\subsection{Etapa 1 --- Teoría de no-unicidad}

\begin{itemize}
\item Redacción formal de $J_4$ (ya está en
  \texttt{tesis/nota\_problema\_inverso.pdf}).
\item Construcción explícita de la familia $\{R:J_4[R]=c\}$
  (cilindro, NC2, ensanche basal, constricción).
\item Webster--Borg si entra.
\item Pregunta bifásica: buscar dos geometrías con el mismo
  $J_4$ (o el mismo factor de escala del recorte) y comparar
  $(q_{\mathrm{ex}},q_{\mathrm{ef}})$ y los perfiles
  $\mathbf{y}(z)$. Si coinciden, la degeneración sobrevivió. Si
  no, medir en qué norma se separan y qué observable la ve
  (candidato: $\phi(0)$, $u_g-u_m$, $z_f$).
\end{itemize}

\subsection{Etapa 2 --- Identificabilidad}

Barrido de parámetros alrededor de Calbuco y de Merapi / Caulle
(las dos ramas). Para cada subconjunto de datos, perfiles de
verosimilitud y SVD de $D\F$. Ruido:
$\sigma_{\log\MER}=\ln 2/2$ (factor~2), más un error de
clasificación de estilo (probabilidad de etiquetar mal una rama
fronteriza). Entregable: una tabla ``parámetro $\times$ dato
$\to$ identificado / acotado / invisible''.

\subsection{Etapa 3 --- Inversión de Calbuco}

Previas $\pi_0$ de la petrología. Verosimilitud sobre
$(\MER_1,\MER_2,\mathrm{masa},H_{\mathrm{columna}},\mathrm{estilo})$.
Muestreo: MCMC si $\F$ es barato (o emulador); si no, algoritmo
de vecindad de Sambridge + remuestreo, como Wong et al.\ (2020).
Verificación predictiva posterior. Comparación con las cotas
independientes de Arzilli / Morgado / Delgado.

\subsection{Etapa 4 --- Geometría como incógnita funcional}

Sintéticos: generar un $d$ con una geometría NC2, invertir con
Tikhonov / TV, mostrar que el mínimo no es único y que el
regularizador elige un representante (el más suave, o el de
menor TV). Reales: si hay evidencia de geometría (cráter,
líticos), usarla como previa, no como ``verdad''.

\subsection{Herramientas}

Python, \texttt{numpy}/\texttt{scipy}, IDA vía
\texttt{scikits.odes}, emulador si hace falta (proceso gaussiano).
El prototipo ya corre sin IDA. No se introduce un framework nuevo
salvo que el muestreo lo exija.

% ===========================================================================
\section{Catálogo de problemas (de dónde se elige el núcleo)}
\label{sec:catalogo}

Las ideas marcadas con $\star$ salieron de conversar el lado
geológico. Dificultad $=$ infraestructura, no calendario.

\paragraph{Familia A --- Parámetros (dimensión finita).}
A1$^\star$ $\Delta P$ (y con ella $P(z)$); A2$^\star$ $c_0$ y
$\xi$; A3 $z_f$ y $\phi_{\mathrm{crit}}$; A4 constantes de
$k(\phi)$. Precisión sobre A1: $P(z)$ no es un dato extra, la
produce el directo. Lo defendible es
\textit{``qué cota superior sobre $\Delta P$ imponen los
observables, y qué dato la acotaría por abajo''}.

\paragraph{Familia B --- Funciones.}
B1$^{\star\star}$ $R(z)$: \textbf{núcleo}. B2 coeficiente de
Navier $\beta(z)$ en la pared (Robin; teoría de unicidad y
estabilidad logarítmica: Chaabane--Jaoua, Inglese, Sincich);
extensión natural del 2D. B3 $\phi(z,0)$ desde una serie de
porosidad de piroclastos. B4 $k(\phi)$ como función.

\paragraph{Familia C --- Historias.}
C1 $P_{\mathrm{cam}}(t)$ desde $\MER(t)$: problema lateral
parabólico, mal puesto clásico (Eldén--Berntsson--Regińska). C2
recarga $Q_{\mathrm{in}}(t)$. C3 tasa de descompresión desde MND
(balance poblacional): conexión directa con petrología de
terreno.

\paragraph{Familia D --- Teoría.}
D1$^\star$ $J_4$ (hecho). D2$^\star$ Webster--Borg. D3
multiplicidad sigmoidal como segunda fuente.

\paragraph{Familia E --- Método.}
E1$^\star$ unificar y regularizar (prerrequisito). E2 adjunto.
E3 selección de criterio de fragmentación (factores de Bayes).
E4 desequilibrio y Darcy (Wong: la permeabilidad controla el
flujo). E5 emulador. E6 diseño óptimo de experimentos. E7
geometría compleja (dique $\to$ cilindro, NC2/NC3 de Aravena).

El núcleo recomendado es D1 + pregunta bifásica + A1--A3 sobre
Calbuco + B1 en sintéticos + el estilo como dato. B2, C1 y D2 son
artículos extras si alcanzan.

% ===========================================================================
\section{Estructura de la tesis (capítulos)}
\label{sec:capitulos}

Pensada para que cada capítulo tenga una frase geológica y una
frase matemática.

\paragraph{Cap.\ 1 --- El volcán y el operador.}
Geología: conducto, estilo, Calbuco, el paper 2025.
Matemática: el DAE, los dos tiros, $\F$ como~\eqref{eq:Fpar}.
Entregable: el operador único (E1).

\paragraph{Cap.\ 2 --- No-unicidad geométrica.}
Geología: seis conductos, $+317\pc$ de volumen, el mismo caudal;
el profundo no se ve.
Matemática: teorema $J_4$, nulo de codimensión~1, SVD.
Entregable: el teorema y las figuras del prototipo, pasados en
limpio.

\paragraph{Cap.\ 3 --- ¿Sobrevive en el bifásico? Estilo.}
Geología: Merapi, Caulle, las dos ramas; el estilo como
observable.
Matemática: contraejemplo o unicidad; el par
$(q_{\mathrm{ex}},q_{\mathrm{ef}})$ como segunda degeneración.
Entregable: la respuesta a la pregunta central.

\paragraph{Cap.\ 4 --- Identificabilidad y el tiempo.}
Geología: qué se gana con series (Van Eaton; fases de Calbuco);
cuándo hace falta el transiente (Deborah).
Matemática: Fisher, perfiles, comparación 1D / transiente / 2D.
Entregable: la tabla parámetro $\times$ dato.

\paragraph{Cap.\ 5 --- Inversión de Calbuco 2015.}
Geología: posteriores sobre $\Delta P$, $c_0$, $z_f$; gatillo;
lectura de depósitos y deformación.
Matemática:~\eqref{eq:bayes}, muestreo, predictiva posterior.
Entregable: el caso real.

\paragraph{Cap.\ 6 --- Geometría como incógnita funcional.}
Geología: NC2/NC3, ensanchamiento sineruptivo (Aravena).
Matemática: Tikhonov / TV, no-unicidad numérica del mínimo.
Entregable: sintéticos +, si hay previa geométrica, Calbuco.

\paragraph{Cap.\ 7 --- Conclusiones para las dos disciplinas.}
Qué puede afirmar un volcanólogo; qué teorema queda; qué se
mediraría mañana.

Si alcanza: apéndice o capítulo corto de Webster (D2) o de Robin
(B2).

% ===========================================================================
\section{Plan de trabajo}
\label{sec:plan}

Sin fechas de calendario. Orden por dependencia; cada etapa tiene
un entregable que se puede mostrar.

\begin{center}
\small
\begin{tabular}{@{}clp{0.62\textwidth}@{}}
\toprule
 & etapa & entregable \\
\midrule
0 & consolidación &
  un $\F$, tests de regresión, fragmentación regularizada,
  decisión sobre inconsistencias \\
1 & teoría &
  $J_4$ + nulo + familia de formas; (opcional) Webster \\
2 & bifásico y estilo &
  sí/no sobrevive $J_4$; mapa de ramas en los casos 2025 \\
3 & identificabilidad &
  tabla Fisher / verosimilitud; Deborah de Calbuco vs.\ un domo \\
4 & Calbuco &
  posteriores + comparación con petrología + (si alcanza)
  factores de Bayes del gatillo \\
5 & $R(\cdot)$ &
  inversión regularizada en sintéticos; efecto del
  regularizador \\
6 & escritura &
  manuscrito; un artículo candidato (caps.\ 2--3: la
  no-unicidad que la volcanología no ha formalizado) \\
\bottomrule
\end{tabular}
\end{center}

Los capítulos 2 y 3 ya tienen resultados preliminares. 4 y 5 son
incrementales. Si hay que recortar, se recorta 6 antes que 3: el
núcleo es la identificabilidad, no la geometría más bonita.

% ===========================================================================
\section{Productos esperados}

\begin{itemize}
\item Tesis de doble titulación, con los siete capítulos de
  \S\ref{sec:capitulos}.
\item Un artículo (caps.\ 2--3): degeneración geométrica +
  pregunta bifásica + estilo como segunda no-unicidad.
\item Un artículo (cap.\ 5), si Calbuco cierra: inversión
  bayesiana de una sub-Pliniana chilena.
\item Código: el operador único, el notebook de ambas ramas, el
  prototipo reproducible (ver \texttt{tesis/prototipo}).
\item Una nota corta para monitoreo (E6), si el diseño óptimo
  sale limpio: qué conviene medir.
\end{itemize}

% ===========================================================================
\section{Riesgos}

\begin{center}
\small
\begin{tabular}{@{}p{0.42\textwidth}p{0.50\textwidth}@{}}
\toprule
riesgo & mitigación \\
\midrule
$\F$ es caro y MCMC no converge &
  emulador (E5) o vecindad de Sambridge, como Wong 2020 \\
\texttt{scikits.odes} no compila &
  el prototipo es numpy/scipy; portar el DAE a un solver
  puro si hace falta \\
Calbuco no restringe el modelo &
  \textbf{ése es el resultado}: se reporta como
  identificabilidad y se complementa con sintéticos \\
los switches impiden gradientes &
  regularizar (E1); si falla, métodos sin derivadas \\
el alcance crece &
  el núcleo son caps.\ 2--3; 4 y 5 son incrementales \\
un profesor lee ``esto es de la otra carrera'' &
  cada capítulo tiene frase G y frase M (\S\ref{sec:capitulos});
  los objetivos están partidos (\S\ref{sec:objetivos}) \\
el recorte se confunde con el modelo &
  ya está escrito en la nota y aquí: se invierte el bifásico;
  el recorte sólo demuestra \\
\bottomrule
\end{tabular}
\end{center}

% ===========================================================================
\section{Preguntas para la reunión}
\label{sec:preguntas}

\paragraph{Para el profesor de matemáticas.}
\begin{enumerate}
\item ¿Determinista (Tikhonov, tasas bajo condición de fuente) o
  bayesiano (posterior, MCMC)? El cap.\ 6 admite ambos; cambia
  el sabor.
\item ¿Adjunto (E2) o diferencias finitas y pocos parámetros?
\item ¿Le basta $J_4$ como teorema central, o quiere unicidad
  condicional (bajo qué datos $R(\cdot)$ \emph{sí} queda)?
\item Webster--Borg: ¿capítulo o sección?
\end{enumerate}

\paragraph{Para el profesor de geología.}
\begin{enumerate}
\item ¿El objetivo es $\Delta P$ en la cámara o el perfil
  $P(z)$? (Véase A1.)
\item ¿Hay serie de porosidad / densidad de piroclastos por
  unidad estratigráfica en Calbuco? Si la hay, B3 es el inverso
  mejor condicionado del catálogo.
\item ¿Hay permeabilidad medida en los clastos? Habilita A4/B4;
  la literatura dice que es el control de primer orden.
\item ¿Qué evidencia independiente hay de la geometría del
  conducto (líticos, cráter, diques)? Previa del cap.\ 6.
\item ¿Confirmamos $T$ y $c_0$? (Advertencia de
  \S\ref{sec:geo}.)
\item Gatillo: ¿le interesa discriminar Arzilli vs.\ Morgado
  como resultado, o lo deja fuera?
\end{enumerate}

\paragraph{Para ambos.}
\begin{enumerate}
\item ¿El caso es Calbuco, o conviene uno con serie de tiempo
  larga (un domo)? La identificabilidad dice que las series
  largas informan más. Es una decisión técnica, no de gusto.
\item ¿Queremos un caso cuasi-estacionario o uno genuinamente
  transiente? Determina qué es invertible
  (\S\ref{sec:math}). Calbuco es lo primero en hidráulica y la
  frontera en cristalización. Un domo cíclico es lo segundo
  por construcción.
\item ¿El estilo (las dos ramas) entra como capítulo propio o
  como sección del operador?
\end{enumerate}

% ===========================================================================
\section{Documentos que acompañan esta propuesta}

\begin{center}
\small
\begin{tabular}{@{}ll@{}}
\toprule
archivo & rol \\
\midrule
\texttt{tesis/propuesta\_tesis\_doble.pdf} &
  \textbf{este documento} (la propuesta formal) \\
\texttt{tesis/PROPUESTA.md} &
  versión de trabajo, catálogo A--E más largo \\
\texttt{tesis/nota\_problema\_inverso.pdf} &
  4~pp.: el inverso sobre el bifásico, $J_4$ en el recorte \\
\texttt{tesis/matematica\_modelo\_bifasico.pdf} &
  8~pp.: el DAE completo, los dos tiros \\
\texttt{tesis/ESTADO\_DEL\_ARTE.md} &
  revisión anotada, DOI verificados \\
\texttt{tesis/presentacion/propuesta\_tesis.pdf} &
  Beamer para la reunión \\
\texttt{tesis/prototipo/} &
  recorte + experimentos 1--4, reproducible \\
\texttt{mainconduit\_ambos.ipynb} &
  laboratorio de las dos ramas \\
\bottomrule
\end{tabular}
\end{center}

\vspace{0.6em}
Reproducir los resultados preliminares:
\begin{verbatim}
pip install numpy scipy matplotlib
cd tesis/prototipo
python3 exp1_no_unicidad_geometria.py
python3 exp2_el_tiempo_rompe_la_degeneracion.py
python3 exp3_identificabilidad.py
python3 exp4_cuando_hace_falta_el_transiente.py
\end{verbatim}

% ===========================================================================
\begin{thebibliography}{99}\small

\bibitem{anderson2013} Anderson, K.\ \& Segall, P.\ (2013).
  \emph{Bayesian inversion of data from effusive volcanic eruptions
  using physics-based models: Application to Mount St.\ Helens
  2004--2008}. JGR 118:2017--2037.
  \href{https://doi.org/10.1002/jgrb.50169}{10.1002/jgrb.50169}

\bibitem{aravena2017} Aravena, A.\ \emph{et al.}\ (2017).
  \emph{Stability of volcanic conduits during explosive eruptions}.
  JVGR.
  \href{https://doi.org/10.1016/j.jvolgeores.2017.03.023}{10.1016/j.jvolgeores.2017.03.023}

\bibitem{aravena2018} Aravena, A.\ \emph{et al.}\ (2018).
  \emph{Conduit stability effects on intensity and steadiness of
  explosive eruptions}. Sci.\ Rep.\ 8.
  \href{https://doi.org/10.1038/s41598-018-22539-8}{10.1038/s41598-018-22539-8}

\bibitem{arzilli2019} Arzilli, F.\ \emph{et al.}\ (2019).
  \emph{The unexpected explosive sub-Plinian eruption of Calbuco
  volcano}. JVGR 378:35--50.
  \href{https://doi.org/10.1016/j.jvolgeores.2019.04.006}{10.1016/j.jvolgeores.2019.04.006}

\bibitem{barmin2002} Barmin, A., Melnik, O.\ \& Sparks, R.S.J.\ (2002).
  \emph{Periodic behavior in lava dome eruptions}. EPSL 199:173--184.

\bibitem{borg1946} Borg, G.\ (1946).
  \emph{Eine Umkehrung der Sturm--Liouvilleschen Eigenwertaufgabe}.
  Acta Math.\ 78:1--96.

\bibitem{castruccio2016} Castruccio, A.\ \emph{et al.}\ (2016).
  \emph{Eruptive parameters and dynamics of the April 2015
  sub-Plinian eruptions of Calbuco volcano}. Bull.\ Volcanol.\ 78(9).
  \href{https://doi.org/10.1007/s00445-016-1058-8}{10.1007/s00445-016-1058-8}

\bibitem{castruccio2025} Castruccio, A., Rebolledo, A.\ \& Gómez, I.\ (2025).
  \emph{The influence of melt composition, temperature, crystallinity
  and water content on eruptive style and eruption rate}.
  JGR Solid Earth 130, e2024JB030599.
  \href{https://doi.org/10.1029/2024JB030599}{10.1029/2024JB030599}

\bibitem{costa2007} Costa, A., Melnik, O.\ \& Sparks, R.S.J.\ (2007).
  \emph{Controls of conduit geometry and wallrock elasticity on lava
  dome eruptions}. EPSL 260:137--151.

\bibitem{delgado2017} Delgado, F.\ \emph{et al.}\ (2017).
  deformación InSAR de Calbuco 2015.

\bibitem{dmv2008} de' Michieli Vitturi, M.\ \emph{et al.}\ (2008).
  \emph{Effects of conduit geometry on magma ascent dynamics}.
  EPSL 272:567--578.

\bibitem{giordano2008} Giordano, D., Russell, J.K.\ \& Dingwell, D.B.\ (2008).
  \emph{Viscosity of magmatic liquids: a model}. EPSL 271:123--134.

\bibitem{hadamard1923} Hadamard, J.\ (1923).
  \emph{Lectures on Cauchy's Problem in Linear Partial Differential
  Equations}. Yale University Press.

\bibitem{kozono2009} Kozono, T.\ \& Koyaguchi, T.\ (2009).
  \emph{Effects of relative motion between gas and liquid on
  1-dimensional steady flow in silicic volcanic conduits}.
  JVGR 180:37--48.
  \href{https://doi.org/10.1016/j.jvolgeores.2008.11.006}{10.1016/j.jvolgeores.2008.11.006}

\bibitem{kozono2012} Kozono, T.\ \& Koyaguchi, T.\ (2012).
  \emph{Effects of gas escape and crystallization on the complexity
  of conduit flow dynamics}. JGR.
  \href{https://doi.org/10.1029/2012JB009343}{10.1029/2012JB009343}

\bibitem{laspina2017} La Spina, G.\ \emph{et al.}\ (2017).
  \emph{Transient numerical model of magma ascent dynamics}.
  JVGR 336:118--139.

\bibitem{massol2001} Massol, H., Jaupart, C.\ \& Pepper, T.W.\ (2001).
  \emph{Ascent and decompression of viscous vesicular magma}.
  JGR 106:16223--16240.

\bibitem{melnik1999} Melnik, O.\ \& Sparks, R.S.J.\ (1999).
  \emph{Nonlinear dynamics of lava dome extrusion}. Nature 402:37--41.

\bibitem{morgado2019} Morgado, E.\ \emph{et al.}\ (2019).
  petrología de Calbuco 2015 (hipótesis de inyección).

\bibitem{raue2009} Raue, A.\ \emph{et al.}\ (2009).
  \emph{Structural and practical identifiability analysis \ldots
  by exploiting the profile likelihood}. Bioinformatics 25:1923--1929.
  \href{https://doi.org/10.1093/bioinformatics/btp358}{10.1093/bioinformatics/btp358}

\bibitem{romero2016} Romero, J.E.\ \emph{et al.}\ (2016).
  \emph{Eruption dynamics of the 22--23 April 2015 Calbuco Volcano}.
  JVGR 317:15--29.

\bibitem{slezin2003} Slezin, Yu.B.\ (2003).
  \emph{The mechanism of volcanic eruptions (a steady state approach)}.
  JVGR 122:7--50.

\bibitem{stuart2010} Stuart, A.M.\ (2010).
  \emph{Inverse problems: a Bayesian perspective}.
  Acta Numerica 19:451--559.
  \href{https://doi.org/10.1017/S0962492910000061}{10.1017/S0962492910000061}

\bibitem{vaneaton2016} Van Eaton, A.R.\ \emph{et al.}\ (2016).
  \emph{Volcanic lightning and plume behavior \ldots Calbuco}.
  GRL.
  \href{https://doi.org/10.1002/2016GL068076}{10.1002/2016GL068076}

\bibitem{wong2017} Wong, Y.Q.\ \emph{et al.}\ (2017).
  \emph{Constraining the magmatic system at Mount St.\ Helens
  \ldots including gas escape and crystallization}. JGR 122:7789--7812.
  \href{https://doi.org/10.1002/2017JB014343}{10.1002/2017JB014343}

\bibitem{wong2019} Wong, Y.Q.\ \& Segall, P.\ (2019).
  \emph{Numerical analysis of time-dependent conduit magma flow
  \ldots Mount St.\ Helens 2004--2008}. JGR.
  \href{https://doi.org/10.1029/2019JB017585}{10.1029/2019JB017585}

\bibitem{wong2020} Wong, Y.Q.\ \emph{et al.}\ (2020).
  \emph{Joint inversions of ground deformation, extrusion flux,
  and gas emissions \ldots}. G-Cubed.
  \href{https://doi.org/10.1029/2020GC009343}{10.1029/2020GC009343}

\end{thebibliography}

\end{document}
